%% cover-letter.tex  —  The American Mathematical Monthly
%% Submission portal: Taylor & Francis ScholarOne
%% Author: Ali Elouafiq, SQLI Digital Laboratory, Levallois-Perret, France.
%% Revised 2026-06-12 (review-panel pass): addressed to the current Editor;
%% page/word counts recomputed from the actual compiled PDFs; figure and
%% reference-style statements added; AI-disclosure specifics deferred to the
%% manuscript's Acknowledgments to keep a single authoritative wording.
%% =============================================================================
\documentclass[12pt]{article}
\usepackage[margin=1in]{geometry}
\usepackage[T1]{fontenc}
\usepackage{times}
\usepackage{color}
\usepackage[normalem]{ulem}
\usepackage{hyperref}
\pagestyle{empty}

\newcommand{\ReviewAdd}[1]{\textcolor{blue}{#1}}
\newcommand{\ReviewDel}[1]{\textcolor{red}{\sout{#1}}}

\begin{document}

\noindent
Ali Elouafiq\\
Principal Research Engineer\\
SQLI Digital Laboratory\\
2--10 rue Thierry Le Luron\\
92300 Levallois-Perret, France\\
\href{mailto:ali@sqli.com}{ali@sqli.com}

\bigskip
\noindent
12 June 2026

\bigskip
\noindent
Professor Annalisa Crannell\\
Editor, \textit{The American Mathematical Monthly}\\
Mathematical Association of America

\bigskip
\noindent
Dear Professor Crannell,

\medskip
\noindent
I submit for your consideration the manuscript:

\begin{center}
  \textbf{The Continuous-Coefficient Jensen Equation:}\\
  \textbf{How Real Mixing Weights Retire Regularity}
\end{center}

\noindent
for publication in \textit{The American Mathematical Monthly} as an
\textbf{Article}.

\bigskip
\noindent\textbf{What the paper does.}
The \ReviewDel{paper}\ReviewAdd{manuscript} shows that the continuous-coefficient Jensen equation
$G(pu_1+(1-p)u_2)=pG(u_1)+(1-p)G(u_2)$ for all $u_1,u_2\in[0,M]$ and
all $p\in[0,1]$ forces $G$ to be affine with \emph{no regularity
hypothesis}: a one-line endpoint substitution ($u_1{:=}M$, $u_2{:=}0$,
$p{:=}v/M$) reads $G(v)$ directly off the straight line through the
endpoints. The \ReviewDel{paper accompanies}\ReviewAdd{manuscript pairs} this one-line observation with
(i)~a dictionary comparing the result with the well-known case of the
discrete-coefficient Jensen equation (where Hamel's pathology is real
and regularity \emph{is} needed), (ii)~a structural explanation of
why the two cases differ, and (iii)~a diagnostic that separates the
saturated-Jensen settings in which these regularity hypotheses are
vestigial from the more common settings in which they remain
load-bearing. The same five hypotheses appear on both sides of the
line: in the axiomatic characterization of Shannon entropy they are
indispensable; in a calibration computation on an atomless
probability space they are not. A worked quantum
example\,---\,mixture-affine assignments on density
operators\,---\,shows the asymmetry \ReviewDel{alive}\ReviewAdd{at work} in a single modern
literature: Gleason-type derivations of the Born rule spend exactly
the dictionary's hypotheses on the measurement side, while the
preparation side, by the convex-domain form of the theorem, spends
none.

\bigskip
\noindent\textbf{Why the Monthly.}
The core result is elementary enough for a second-year analysis student
to verify in an afternoon, \ReviewDel{and yet}\ReviewAdd{yet} modern applied work spans both
sides of a sharp line: in some saturated-Jensen settings the
continuous-coefficient equation makes the classical regularity
hypotheses vestigial, while in others\,---\,superficially
identical\,---\,those hypotheses remain indispensable. The article
supplies the diagnostic that tells them apart, and \ReviewDel{(in its closing
note on provenance) recounts}\ReviewAdd{its closing note on provenance recounts} how the author was led to it by an
instance of the rarer error of adding a hypothesis that was not
needed. \ReviewDel{The Monthly is the natural home for}\ReviewAdd{This is a natural fit for the \textit{Monthly}:} an expository article that
makes this asymmetry explicit, mirrors the regularity dictionary, and
invites readers to audit their own derivations.

\bigskip
\noindent\textbf{Novelty and relationship to prior art.}
The substitution technique that closes Theorem~1 is folklore:
Acz\'{e}l's 1966 monograph (\S2.1.4) contains a dyadic precursor proved
under a continuity hypothesis\ReviewDel{, and the genre is what Kuczma's 2009
introduction describes as the part of Jensen-equation theory that does
not require regularity assumptions on the function}. The precise
continuous-coefficient regularity-free statement of Theorem~1 does
not, however, appear in this exact form in the standard references the
author consulted. The contribution of this article is therefore not
Theorem~1 itself but its \ReviewDel{\emph{explicit packaging}}\ReviewAdd{explicit framing}: the dictionary
contrasting the two Jensen forms, the structural explanation of the
Hamel-pathology mechanism, and the identification of a recurrence
pattern in derivations that produce the continuous-coefficient form as
a saturated identity. \ReviewDel{A closing note records}\ReviewAdd{The closing note records} that the observation
reached the author through a formalization project on a separate
topic\,---\,a route that may interest readers curious about how
mechanized proof can pressure-test classical heuristics.

\bigskip
\noindent\textbf{Submission details.}
\begin{itemize}\setlength\itemsep{2pt}
  \item \textbf{Manuscript type:} Article (expository).
  \item \textbf{Page count:} 19 pages anonymous body; 20 pages with the
        title page (Monthly \LaTeX{} template formatting).
  \item \textbf{Word count (approximate):} 7{,}400 words of main text,
        excluding figures, tables, and references.
  \item \textbf{Figures and tables:} three monochrome figures (drawn in
        Ti\textit{k}Z, embedded in the \LaTeX{} source per the
        \emph{Monthly}'s figure instructions) and two tables.
  \item \textbf{References:} 54 entries, typeset in the \emph{Monthly}'s
        house style (alphabetical, bracketed numbers, DOIs included
        where available).
  \item \textbf{MSC 2020 primary:} 39B22 (Functional equations for real functions).
  \item \textbf{MSC 2020 secondary:} 26B25 (Convexity), 91B16 (Utility theory),
        94A17 (Measures of information; entropy).
  \item \textbf{Double-anonymous:} Yes. The anonymous body
        (\texttt{manuscript-anon.pdf}) contains no author-identifying
        information; the file with author details
        (\texttt{manuscript.pdf}) carries the title page on page~1.
  \item \textbf{Preprint status:} The manuscript has not been deposited on
        arXiv. It is not under consideration elsewhere.
  \item \textbf{Competing interests:} None.
  \item \textbf{Funding:} This research received no specific grant from any
        funding agency in the public, commercial, or not-for-profit
        sectors.
  \item \textbf{Use of generative AI:} Disclosed in the manuscript's
        Acknowledgments section (AI-assisted exploratory
        theorem-discovery within a Lean~4 formalization harness, and
        AI-assisted proofreading and formatting). All results and their
        presentation were analyzed and carried out by the author, who
        assumes full responsibility for correctness, originality, and
        content.
\end{itemize}

\bigskip
\noindent
I look forward to the referees' comments.

\bigskip
\noindent
Sincerely,

\vspace{28pt}
\noindent
Ali Elouafiq
\end{document}
